\clearpage
\fancyhead[L]{\small OBSERVING DEFLECTIONS.}
\SourceLeaf{401}
\pdfbookmark[0]{On a way to make deflections easier to observe. 1839}{ablenkungen-1839}
\begin{center}
{\large ON A WAY\par\medskip TO MAKE DEFLECTIONS EASIER TO OBSERVE.}
\par\medskip
Resultate aus den Beobachtungen des magnetischen Vereins. 1839. II.
\end{center}
\par\medskip\Needspace{4\baselineskip}\pdfbookmark[1]{1.}{ablenkungen-art-1}\centerline{\textbf{1.}}\nobreak\smallskip
If, in addition to the Earth’s magnetic force, another force acts steadily on the needle of a magnetometer in a direction inclined to the magnetic meridian, the needle takes up a new equilibrium position. The magnitude of the deflection can serve to measure the additional force. For the deflection to be measurable, the new equilibrium position must still lie within the scale. And if one cannot or does not wish to wait for the needle to come completely to rest, the remaining swings must also stay within the limits of the scale. Suppose the needle was at rest while the Earth’s magnetism alone acted on it, and the additional force is then brought into full effect all at once. The needle begins a swing whose midpoint is the new equilibrium position. Its former position is one point of elongation, and the other falls just as far from the midpoint on the opposite side. If the new equilibrium position is within the scale but close to its limit, then with the arc of swing decreasing slowly, and without artificial aids, one would have to wait a long time before that point could be determined in this way. In every case, however, the hourly change in declination alone would greatly diminish the reliability and usefulness of the determination, and the determination would be rendered almost entirely useless
\expandafter\def\csname HistoricalNoteText404\endcsname{Resultate, 1837. IV. [p. 383 of this volume.]}
\renewcommand{\Article}[1]{\par\medskip\Needspace{4\baselineskip}\pdfbookmark[1]{#1.}{ablenkungen-art-#1}\centerline{\textbf{#1.}}\nobreak\smallskip}%
\SourceLeaf{402}in cases where the strength of the additional force changes considerably even in a short time, as with galvanic currents.

\Article{2}
The following simple procedure remedies this difficulty. First let the additional force act for only a third of the swing duration. Then suspend it for an interval of the same length, and only then let it act continuously. For example, suppose the swing duration of the magnetometer needle is 30 seconds and we wish to measure the deflection produced by a galvanic current. Close the circuit at a beat of the seconds clock, which is counted as 0, open it again at $10''$, and close it again at $20''$ and leave it closed. If the deflection is to be produced by a magnetic bar laid at a particular spot, approach the spot, which has been marked beforehand precisely and conveniently, with the bar initially held vertically. Lay it down abruptly at $0''$, stand it up again just as quickly at $10''$, and lay it down a second time at $20''$ and leave it there. The result is that for the first 10 seconds the needle approaches the position corresponding to the deflection from its original position of rest with increasing speed. At $10''$ it has reached exactly the midpoint between the two positions. During the second 10 seconds it covers the other half of the distance with slowing motion, so that at 20 seconds it will have reached the new position and lost all motion.

It is easy to see that in much the same way the needle can be moved from one stationary deflected position to the opposite one so that it is motionless there. Let the deflecting force act in the opposite direction for a third of the swing duration, then in its previous direction again for just as long, and then reverse it once more. With galvanic currents a suitable commutator makes the reversal nearly instantaneous. With deflecting magnetic bars a quick half-turn, most conveniently a horizontal one, brings the north pole of the bar to the place of the south pole.

Finally, once the deflection has been observed, it is clear that the needle can be brought to rest again in the undisturbed magnetic meridian in the same way. One need only suspend the deflecting force first for a third of the swing du\SourceLeaf{403}ration, then let it act again for just as long, and finally let it cease.

\Article{3}
The procedure described rests on the assumption that

\textit{first}, the needle swings in such a way that its distance from the midpoint of the swing, as long as that midpoint is itself unaltered, remains proportional to the sine of an angle that changes uniformly and increases by $180^0$ during one swing duration, and

\textit{second}, the additional force leaves the swing duration unchanged.

Insofar as the two assumptions fail to hold \textit{with absolute precision}, and in carrying out the procedure the reversals can neither be quite instantaneous nor the prescribed intervals kept with absolute accuracy, the needle will indeed seldom be found perfectly at rest at the end of the operation. For practical purposes, however, it is enough if the remaining motion is so slight that one can at once begin to observe the true equilibrium position in the usual way.

Under the circumstances in question, these assumptions can depart only very slightly from the truth. The applicability of the magnetometer itself depends on the additional force producing only a moderate deflection. With such a deflection, apart from an exception to be mentioned shortly, the law in the first assumption holds with sufficient accuracy. The change in swing duration caused by the deflecting force is quite imperceptible if the force acts perpendicular to the magnetic meridian, as it almost always does. Even if it acted in an oblique direction, however, provided that it was only a small fraction of the Earth’s magnetic force, the resulting change in swing duration would still be quite immaterial over the short time of the operation.

There is just one exception, namely when the needle swings under the influence of a damper that substantially reduces the size of the arc of swing. In that case the law above no longer holds, and following the procedure described above exactly would fail to achieve the desired result. On the other hand, the
\SourceLeaf{404}difficulty noted in Article 1 is then much less serious, since a powerful damper brings the needle to rest of itself in a moderate time. Yet in this case the procedure requires only a modification to achieve the same result. Since it is always desirable to avoid any unnecessary loss of time, it is worthwhile from both a practical and a theoretical point of view. to consider the question in full generality.

\Article{4}
We must first solve the following general problem.

A magnetic bar swings while the forces acting on it are repeatedly changed. The swing duration and the logarithmic decrement\HistoricalNote{404} nevertheless remain unchanged, and the arcs of swing remain small enough for quantities of the third order to be neglected. From the initial state of motion we are to deduce the state after the final change.

Let $T$ be the swing duration, $\epsilon$ the logarithmic decrement, $e$ the base of the hyperbolic logarithms, $m$ the modulus of the Briggs logarithms, and $\pi$ the ratio of the circumference of a circle to its diameter. Set
\[
n=\frac{\pi}{T},\qquad \epsilon=\frac{\lambda}{mT}
\]
Under the assumptions above, the position $x$ at time $t$ is accordingly given by the formula
\[
x=p+Ae^{-\epsilon t}\sin(nt-B)
\]
which can also be written in the form
\[
x=p+ae^{-\epsilon t}\cos nt+be^{-\epsilon t}\sin nt
\]
where $p$ denotes the equilibrium position, and the coefficients $a,b$ remain constant as long as $p$ is constant. This gives the velocity as
\[
\frac{dx}{dt}=-e^{-\epsilon t}(na\sin nt+\epsilon a\cos nt-nb\cos nt+\epsilon b\sin nt)
\]
or, if we introduce an auxiliary angle $\varphi$ such that $\frac{\epsilon}{n}=\operatorname{tang}\varphi$,
\SourceLeaf{405}\[
\frac{dx}{dt}=-\frac{ne^{-\epsilon t}}{\cos\varphi}(a\sin(nt+\varphi)-b\cos(nt+\varphi))
\]
We write $ae^{-\epsilon t}\cos nt+be^{-\epsilon t}\sin nt$ as $u$, so that $x=p+u$.

Now let $t',t'',t'''$ be the particular values of $t$ at which the acting force is changed. Let the particular values of $p,a,b$ in the different time intervals be as follows
\[
\begin{array}{lll}
p^0,\ a^0,\ b^0&\text{before}&t'\\
p',\ a',\ b'&\text{from}&t'\ \text{to}\ t''\\
p'',\ a'',\ b''&\text{from}&t''\ \text{to}\ t'''\\
p''',\ a''',\ b'''&\text{after}&t'''
\end{array}
\]
Finally, let the general expression for $u$, when the particular values of $a$ and $b$ are substituted, become $u^0,u',u'',u'''$, so that before the first change $x=p^0+u^0$, from then until the second $x=p'+u'$, and so on.

Since the instant $t'$ is both the last of the first time interval and the first of the following one, at $t=t'$ both $x$ and $\frac{dx}{dt}$ must take the same values whether we substitute for $p,a,b$ in the general expressions above the values $p^0,a^0,b^0$ or $p',a',b'$.

Thus
\[
\begin{aligned}
0&=p'-p^0+(a'-a^0)e^{-\epsilon t'}\cos nt'+(b'-b^0)e^{-\epsilon t'}\sin nt'\\
0&=(a'-a^0)\sin(nt'+\varphi)-(b'-b^0)\cos(nt'+\varphi)
\end{aligned}
\]
from which we easily deduce
\[
\begin{aligned}
a'-a^0&=-\frac{p'-p^0}{\cos\varphi}.e^{\epsilon t'}\cos(nt'+\varphi)\\
b'-b^0&=-\frac{p'-p^0}{\cos\varphi}.e^{\epsilon t'}\sin(nt'+\varphi)
\end{aligned}
\]
and hence
\[
u'=u^0-\frac{p'-p^0}{\cos\varphi}.e^{-\epsilon(t-t')}\cos(n(t-t')-\varphi)
\]
In the same way we obtain
\[
\begin{aligned}
u''&=u'-\frac{p''-p'}{\cos\varphi}.e^{-\epsilon(t-t'')}\cos(n(t-t'')-\varphi)\\
u'''&=u''-\frac{p'''-p''}{\cos\varphi}.e^{-\epsilon(t-t''')}\cos(n(t-t''')-\varphi)
\end{aligned}
\]
and so on if further changes in the forces producing the motion occur.

\SourceLeaf{406}The initial state of motion thus determines each subsequent state.

\Article{5}
In the case under investigation we must set $p''=p^0$ and $p'''=p'$. This gives
\[
\begin{aligned}
u'''={}&u^0-\frac{p'-p^0}{\cos\varphi}.e^{-\epsilon t}.\\
&.\{e^{\epsilon t'}\cos(n(t-t')-\varphi)-e^{\epsilon t''}\cos(n(t-t'')-\varphi)\\
&\qquad+e^{\epsilon t'''}\cos(n(t-t''')-\varphi)\}
\end{aligned}
\]
If we set
\[
\begin{aligned}
e^{-\epsilon(t''-t')}\cos n(t''-t')-1+e^{\epsilon(t'''-t'')}\cos n(t'''-t'')&=f\\
e^{-\epsilon(t''-t')}\sin n(t''-t')\phantom{-1}{}-e^{\epsilon(t'''-t'')}\sin n(t'''-t'')&=g
\end{aligned}
\]
the formula becomes
\[
\begin{aligned}
u'''={}&u^0-\frac{p'-p^0}{\cos\varphi}.e^{-\epsilon(t-t'')}.\{f\cos(n(t-t'')-\varphi)\\
&\qquad-g\sin(n(t-t'')-\varphi)\}
\end{aligned}
\]
It follows that if the intervals $t''-t',t'''-t''$ are chosen so that $f=0$ and $g=0$, then in general
\[
u'''=u^0
\]
or $a'''=a^0,b'''=b^0$.

Thus if the needle was at rest at $p^0$ before the changes, it will be at rest at $p'$ after them. In the opposite case, at every instant after the three changes the needle will have exactly the same velocity and the same position relative to the centre of its motion $p'$ that it would have relative to $p^0$ at the same instant if it had continued its original motion undisturbed. In short, only the centre of the motion is displaced, and the motion itself is entirely unchanged.

\Article{6}
It remains to determine the intervals so that the equations $f=0,g=0$ are satisfied. If we set
\SourceLeaf{407}\[
t''-t'=qT,\qquad t'''-t''=rT
\]
and recall that $e=10^m$, those equations become
\[
\begin{aligned}
10^{-q\lambda}\cos q\pi+10^{r\lambda}\cos r\pi&=1\\
10^{-q\lambda}\sin q\pi&=10^{r\lambda}\sin r\pi
\end{aligned}
\]
If the arc of swing decreases imperceptibly, we must therefore set $\cos q\pi+\cos r\pi=1$ and $\sin q\pi=\sin r\pi$, and hence
\[
q\pi=r\pi=60^0\ \text{or}\ =\tfrac{1}{3}\pi,\quad\text{and}\quad t''-t'=t'''-t''=\tfrac{1}{3}T
\]
as already noted in Article 2. With an appreciable logarithmic decrement, however, those equations will have to be solved indirectly. The calculation can be put in the following form.

Combining the equations gives
\[
\begin{gathered}
\operatorname{tang}r\pi=\frac{\sin q\pi}{10^{q\lambda}-\cos q\pi},\qquad
10^{2r\lambda}=1-2.10^{-q\lambda}\cos q\pi+10^{-2q\lambda}
\end{gathered}
\]
Eliminating $r$ therefore gives the equation with one unknown quantity
\[
\frac{\pi}{2\lambda}\log(1-2.10^{-q\lambda}\cos q\pi+10^{-2q\lambda})=\operatorname{Arc.\,tang}\frac{\sin q\pi}{10^{q\lambda}-\cos q\pi}
\]
where the Briggs logarithm is meant. Once this equation has been satisfied, we evidently have the value of $r$ as well.

\Article{7}
To spare those who wish to carry out the procedure described with a damper the calculation explained in the preceding article, I give here a table computed in advance, from which the ratio of each of the two intervals to the swing duration can be read off at once for every logarithmic decrement. It shows that as the logarithmic decrement increases, the first interval keeps growing and the second keeps shrinking. Their sum is exactly equal to two thirds of the swing duration only for $\lambda=0$, but departs from it much more slowly. It hardly needs pointing out that in practice it is enough to consider, say, just the first decimal places of the values of $q$ and $r$.

\SourceLeaf{408}\begin{center}\textit{\textbf{Table.}}\end{center}
\begingroup\fontsize{9}{11}\selectfont
\begin{center}
\begin{tabular}{c|c|c||c|c|c}
$\lambda$&$q$&$r$&$\lambda$&$q$&$r$\\\hline
0&0.33333&0.33333&0.30&0.45921&0.21406\\
0.01&0.33757&0.32911&0.31&0.46322&0.21048\\
0.02&0.34181&0.32489&0.32&0.46721&0.20694\\
0.03&0.34606&0.32068&0.33&0.47118&0.20343\\
0.04&0.35031&0.31648&0.34&0.47513&0.19996\\
0.05&0.35456&0.31229&0.35&0.47906&0.19652\\[3pt]
0.06&0.35882&0.30812&0.36&0.48297&0.19311\\
0.07&0.36308&0.30395&0.37&0.48685&0.18975\\
0.08&0.36734&0.29981&0.38&0.49071&0.18641\\
0.09&0.37160&0.29568&0.39&0.49454&0.18311\\
0.10&0.37585&0.29156&0.40&0.49835&0.17985\\[3pt]
0.11&0.38011&0.28746&0.41&0.50214&0.17663\\
0.12&0.38436&0.28338&0.42&0.50590&0.17344\\
0.13&0.38861&0.27932&0.43&0.50963&0.17029\\
0.14&0.39285&0.27528&0.44&0.51334&0.16718\\
0.15&0.39708&0.27126&0.45&0.51702&0.16411\\[3pt]
0.16&0.40131&0.26727&0.46&0.52067&0.16107\\
0.17&0.40552&0.26329&0.47&0.52430&0.15808\\
0.18&0.40973&0.25934&0.48&0.52790&0.15512\\
0.19&0.41393&0.25542&0.49&0.53147&0.15220\\
0.20&0.41812&0.25152&0.50&0.53501&0.14931\\[3pt]
0.21&0.42230&0.24764&0.51&0.53852&0.14647\\
0.22&0.42646&0.24379&0.52&0.54201&0.14367\\
0.23&0.43061&0.23997&0.53&0.54546&0.14090\\
0.24&0.43474&0.23618&0.54&0.54889&0.13817\\
0.25&0.43886&0.23242&0.55&0.55229&0.13548\\[3pt]
0.26&0.44297&0.22868&0.56&0.55566&0.13283\\
0.27&0.44705&0.22498&0.57&0.55900&0.13022\\
0.28&0.45112&0.22131&0.58&0.56231&0.12765\\
0.29&0.45517&0.21767&0.59&0.56559&0.12511\\
0.30&0.45921&0.21406&0.60&0.56884&0.12261
\end{tabular}
\end{center}\endgroup

\Article{8}
The assumption underlying our theory that the three changes occur instantaneously never holds with complete precision when the procedure is actually carried out, although for deflections produced by galvanic currents the time needed for each change may be regarded as imperceptible. For deflections produced by magnetic bars, however, that time is more or less considerable according to their size and weight. With twenty-five-pound bars a change will always take several seconds, especially when it is a change between two opposite posi\SourceLeaf{409}tions rather than between a vertical and a horizontal position. In this case, which is in fact by far the most important and the most common, the operation can easily be arranged so that the result is hardly affected. One need only take care that the second and third changes are made in the same way as the first and thus occupy an equal time, and subtract that time from the intervals that would otherwise be needed. For example, if (as in \textit{Res.} 1837. IV.) [p. 390 of this volume] the logarithmic decrement is 0.33570 and the swing duration is $21''21439$, then the table above gives the first interval $=10''04$ and the second $=4''27$. If three seconds are found necessary to carry out a change, one \textit{begins} the first change at $0''$. From $3''$ to $10''$. the bar remains in the new position. Through the new change \textit{beginning} at $10''$ the bar is back in its first position at $13''$. It stays in that position for only $1\tfrac{1}{4}$ seconds, after which the third change begins, so the whole operation is completed only at $17\tfrac{1}{4}$ seconds. A more extensive investigation, which must be passed over here for lack of space, shows that if $p^0$ changes to $p'$ gradually rather than in a jump, and likewise at the second change $p'$ to $p^0$, and at the third again $p^0$ to $p'$, the result remains exactly as stated at the end of Article 5, provided only that the three transition times are equally long, that the three transitions themselves proceed through similar sequences of stages, and that the calculated intervals $qT,rT$ are reckoned from the instants when the changes begin, or, equivalently, that the first two transition times are included in those intervals.

